% GENERATED FILE. Edit canonical lessons, then run npm run generate:books.
% canonical-source-hash: fnv1a64:050c069e342dfe4e
% canonical-lessons: GE-C376-pauken, GE-C376-auswendig, GE-C376-konzentrieren, GE-C376-konzentration, GE-C376-wiederholung

\chapter{Für die Prüfung lernen}
\label{ch:ge-fur-die-prufung-lernen}
\hlchaptermodality{376}

\begin{chapteropening}
\textbf{By the end of this chapter:} \emph{I can talk about cramming, learning by heart, concentrating and revising for an exam.}

Complete the last lesson of chapter 376: I can talk about cramming, learning by heart, concentrating and revising for an exam.
\end{chapteropening}

\section[pauken]{pauken — to cram, to swot (informal)}

\label{lesson:GE-C376-pauken}



\begin{quote}
Before the new one: say the German for written, then the German for the school report.
\end{quote}

\subsection*{You'll want to know: pauken}
\textbf{pauken} — \textquotedblleft{}to cram, to swot (informal)\textquotedblright{}.

\begin{quote}
Vor der Prüfung hat sie jeden Abend Vokabeln gepaukt.
\end{quote}

\begin{quote}
\emph{Before the exam she crammed vocabulary every evening.}
\end{quote}

Informal; \emph{büffeln} means the same. In Austria people also say \emph{stucken}.

\begin{cousinweb}
\emph{pauken} is built on \emph{Pauke} \textquotedblleft{}kettledrum\textquotedblright{}, from Middle High German \emph{pūke} (EtymDB). The school sense is usually explained as drumming facts into your head.
\end{cousinweb}

\begin{grammarlens}[title={What you've built}]
One more word for telling, explaining and describing at B1.
\end{grammarlens}

\subsection*{Your turn}
\begin{itemize}
\raggedright
  \item \emph{Say it:} \emph{pauken}
  \item \emph{Say it:} \emph{pauken}, once more
  \item \emph{Say it:} a sentence of your own with \emph{pauken}
  \item \emph{Recall:} say the German for written, then the German for the school report, then say \emph{pauken} again
\end{itemize}

\begin{tcolorbox}[breakable,colback=teal!4,colframe=teal!35!black,title={Before you move on}]
What does pauken mean? (\textquotedblleft{}To cram, to swot (informal)\textquotedblright{}.) Say it once more.
\end{tcolorbox}
\section[auswendig]{auswendig — by heart}

\label{lesson:GE-C376-auswendig}



\begin{quote}
Before the new one: say the German for the school report, then the German for to cram.
\end{quote}

\subsection*{You'll want to know: auswendig}
\textbf{auswendig} — \textquotedblleft{}by heart\textquotedblright{}.

\begin{quote}
Die Kinder lernen das Lied auswendig.
\end{quote}

\begin{quote}
\emph{The children are learning the song by heart.}
\end{quote}

\emph{Auswendig lernen} — to learn by heart; \emph{auswendig können} — to know by heart.

\begin{cousinweb}
\emph{auswendig} is built from \emph{aus} \textquotedblleft{}out\textquotedblright{} and \emph{wenden} \textquotedblleft{}to turn\textquotedblright{}, Old High German \emph{wenten} (EtymDB), and first meant \textquotedblleft{}on the outside\textquotedblright{}: you can say the text without looking into the book.
\end{cousinweb}

\begin{grammarlens}[title={What you've built}]
One more word for telling, explaining and describing at B1.
\end{grammarlens}

\subsection*{Your turn}
\begin{itemize}
\raggedright
  \item \emph{Say it:} \emph{auswendig}
  \item \emph{Say it:} \emph{auswendig}, once more
  \item \emph{Say it:} a sentence of your own with \emph{auswendig}
  \item \emph{Recall:} say the German for the school report, then the German for to cram, then say \emph{auswendig} again
\end{itemize}

\begin{tcolorbox}[breakable,colback=teal!4,colframe=teal!35!black,title={Before you move on}]
What does auswendig mean? (\textquotedblleft{}By heart\textquotedblright{}.) Say it once more.
\end{tcolorbox}
\section[konzentrieren]{konzentrieren — to concentrate (mostly reflexive, sich konzentrieren)}

\label{lesson:GE-C376-konzentrieren}



\begin{quote}
Before the new one: say the German for to cram, then the German for by heart.
\end{quote}

\subsection*{You'll want to know: konzentrieren}
\textbf{konzentrieren} — \textquotedblleft{}to concentrate (mostly reflexive, sich konzentrieren)\textquotedblright{}.

\begin{quote}
Bei dem Lärm kann ich mich nicht konzentrieren.
\end{quote}

\begin{quote}
\emph{I cannot concentrate with this noise.}
\end{quote}

Mostly reflexive: \emph{sich auf etwas konzentrieren}. Perfekt: \emph{hat sich konzentriert}.

\begin{cousinweb}
\emph{konzentrieren} was borrowed from French \emph{concentrer} (EtymDB), built from \emph{con-} \textquotedblleft{}together\textquotedblright{} and \emph{centre}, from Latin \emph{centrum} \textquotedblleft{}the middle point\textquotedblright{} (Lewis \& Short): you bring everything to one centre.
\end{cousinweb}

\begin{grammarlens}[title={What you've built}]
One more word for telling, explaining and describing at B1.
\end{grammarlens}

\subsection*{Your turn}
\begin{itemize}
\raggedright
  \item \emph{Say it:} \emph{konzentrieren}
  \item \emph{Say it:} \emph{konzentrieren}, once more
  \item \emph{Say it:} a sentence of your own with \emph{konzentrieren}
  \item \emph{Recall:} say the German for to cram, then the German for by heart, then say \emph{konzentrieren} again
\end{itemize}

\begin{tcolorbox}[breakable,colback=teal!4,colframe=teal!35!black,title={Before you move on}]
What does konzentrieren mean? (\textquotedblleft{}To concentrate (mostly reflexive, sich konzentrieren)\textquotedblright{}.) Say it once more.
\end{tcolorbox}
\section[die Konzentration]{die Konzentration — concentration}

\label{lesson:GE-C376-konzentration}



\begin{quote}
Before the new one: say the German for by heart, then the German for to concentrate.
\end{quote}

\subsection*{You'll want to know: die Konzentration}
\textbf{die Konzentration} — \textquotedblleft{}concentration\textquotedblright{}.

\begin{quote}
Nach zwei Stunden lässt meine Konzentration nach.
\end{quote}

\begin{quote}
\emph{After two hours my concentration fades.}
\end{quote}

No plural in this sense. \emph{Nachlassen} — to ease off, to fade.

\begin{cousinweb}
\emph{Konzentration} goes with \emph{konzentrieren}, from French \emph{concentrer} (EtymDB), built on \emph{centre}, from Latin \emph{centrum} \textquotedblleft{}the middle point\textquotedblright{} (Lewis \& Short).
\end{cousinweb}

\begin{grammarlens}[title={What you've built}]
One more word for telling, explaining and describing at B1.
\end{grammarlens}

\subsection*{Your turn}
\begin{itemize}
\raggedright
  \item \emph{Say it:} \emph{die Konzentration}
  \item \emph{Say it:} \emph{die Konzentration}, once more
  \item \emph{Say it:} a sentence of your own with \emph{die Konzentration}
  \item \emph{Recall:} say the German for by heart, then the German for to concentrate, then say \emph{die Konzentration} again
\end{itemize}

\begin{tcolorbox}[breakable,colback=teal!4,colframe=teal!35!black,title={Before you move on}]
What does die Konzentration mean? (\textquotedblleft{}Concentration\textquotedblright{}.) Say it once more.
\end{tcolorbox}
\section[die Wiederholung]{die Wiederholung — the revision (for an exam); the repetition}

\label{lesson:GE-C376-wiederholung}



\begin{quote}
Before the new one: say the German for to concentrate, then the German for concentration.
\end{quote}

\subsection*{You'll want to know: die Wiederholung}
\textbf{die Wiederholung} — \textquotedblleft{}the revision (for an exam); the repetition\textquotedblright{}.

\begin{quote}
Vor der Klausur machen wir eine Wiederholung des ganzen Stoffes.
\end{quote}

\begin{quote}
\emph{Before the exam we revise all the material.}
\end{quote}

Plural: \emph{die Wiederholungen}. From \emph{wiederholen} — to repeat, to revise. On television, \emph{die Wiederholung} is a repeat.

\begin{cousinweb}
\emph{Wiederholung} is formed from \emph{wiederholen} with \emph{-ung} (EtymDB); the verb joins \emph{wieder} \textquotedblleft{}again\textquotedblright{} to \emph{holen} \textquotedblleft{}to fetch\textquotedblright{}, from Middle High German \emph{holen} (EtymDB): you fetch the material again.
\end{cousinweb}

\begin{grammarlens}[title={What you've built}]
One more word for telling, explaining and describing at B1.
\end{grammarlens}

\subsection*{Your turn}
\begin{itemize}
\raggedright
  \item \emph{Say it:} \emph{die Wiederholung}
  \item \emph{Say it:} \emph{die Wiederholung}, once more
  \item \emph{Say it:} a sentence of your own with \emph{die Wiederholung}
  \item \emph{Recall:} say the German for to concentrate, then the German for concentration, then say \emph{die Wiederholung} again
\end{itemize}

\begin{tcolorbox}[breakable,colback=teal!4,colframe=teal!35!black,title={Before you move on}]
What does die Wiederholung mean? (\textquotedblleft{}The revision (for an exam); the repetition\textquotedblright{}.) Say it once more.
\end{tcolorbox}
